\section[Introduction]{Introduction}

\begin{frame}{Course Content}
	Principal Component Analysis (PCA) is an \textbf{unsupervised learning} algorithm aimed at reducing the \textbf{dimensionality} of a regression or classification problem. In this context, dimensionality is related to the input features $x \in \mathcal{X}$.
	\vspace{0.5cm}
	
	In this course, we will cover the following topics:
	\begin{itemize}
		\pause
		\item Constrained optimization using \textbf{Lagrange multipliers}.
		\pause
		\item The challenges of dimensionality reduction for \textbf{statistical learning}.
		\pause
		\item Properties of \textbf{linear algebra} useful for PCA.
		\pause
		\item Practical application to \textbf{image processing}.
	\end{itemize}
\end{frame}